\label{sec:conclusion}

Algorithmic redistricting produces substantially more competitive
congressional districts than enacted partisan maps: 85 competitive districts
versus 65 under the enacted plans used for the 117th Congress, a 31\%
increase.  Among competitive algorithmic districts, the partisan split is
nearly symmetric (43 Democratic-leaning, 42 Republican-leaning), and the
majority remain durably competitive under vote-share shifts of up to 4
percentage points.

These findings support the democratic-responsiveness argument for algorithmic
redistricting, but with a critical qualification: competitiveness is a
\emph{byproduct} of geometric neutrality, not a design target.  An algorithm
that targeted competitiveness would need to access partisan data, undermining
the anti-gerrymandering rationale for algorithmic redistricting in the first
place.  The appropriate framing is that enacted maps actively suppress
competition by design, and algorithmic maps restore the competition that
underlying geography supports.

Several limitations bound our conclusions.  First, we use a single election
cycle (2020) to measure competitiveness; a more robust analysis would use
average presidential vote over multiple cycles.  Second, we do not account
for incumbency advantage, which can insulate nominally competitive districts
from genuine electoral contestation.  Third, our algorithmic maps use a
single randomization seed; an ensemble analysis \citep{dellaLibera2026uncertainty}
would characterize the range of competitive-district outcomes achievable under
the algorithm's constraints.

The state heterogeneity analysis identifies North Carolina, Texas, Ohio, and
Florida as the states with the largest competitiveness gains from algorithmic
redistricting---precisely the states where enacted maps have attracted the
most sustained legal and political challenges.  This convergence is not
coincidental: the states where gerrymandering is most aggressive are also the
states where algorithmic redistricting would produce the largest electoral
benefits.

The near-symmetric partisan split of competitive algorithmic districts
addresses a concern sometimes raised that algorithmic redistricting
inadvertently advantages one party.  Under the 2020 vote baseline, algorithmic
maps allocate competitive seats almost exactly in proportion to the underlying
partisan balance.  This symmetry is itself evidence that the algorithm is
doing what it is supposed to do: responding to geography rather than to party.

For policymakers and courts considering algorithmic redistricting as a remedy
or structural reform, the competitive-elections finding provides concrete
democratic benefit that is measurable and reproducible.  The 20 additional
competitive seats represent 20 additional districts where incumbents must
maintain genuine constituency contact and where voters retain the credible
power to sanction their representatives through electoral challenge.
